\section[Introdução]{Introdução}

%   Na introdução o aluno deve apresentar os seguintes itens: Descrição do Projeto,
% Stakeholders do Projeto, Período de Execução, Professor Orientador e uma foto do
% time.

A empresa WeConecta conduz este projeto para transformar sua solução atual, integrada ao WhatsApp e dependente de personalizações feitas pelo time interno, em uma plataforma Web com interfaces intuitivas e foco na autonomia do usuário. Na área da saúde, onde a análise de dados impulsiona diagnósticos precoces e melhores resultados clínicos, a proposta é eliminar gargalos de configuração e acelerar a criação, a distribuição e o acompanhamento de questionários personalizados por parte dos profissionais de saúde.

O objetivo central é disponibilizar um ambiente Web em que clientes da área da saúde acessem e personalizem rapidamente formulários criados pela WeConecta, permitindo que pacientes respondam por meio de links públicos e, quando apropriado, por uma interface de chat conversacional para coleta de dados estruturados. Com isso, espera-se agilizar a triagem, aprimorar a prevenção de riscos e reduzir a dependência do time da WeConecta para ajustes simples.

O projeto foi realizado em parceria com a WeConecta Ltda., tendo como principais \textit{stakeholders} os colaboradores João Jacques e Marcelo Nozari. A orientação acadêmica foi conduzida pela Prof.ª Dra. Cristina Moreira Nunes, que acompanhou o escopo, a metodologia e a coerência técnica ao longo das etapas. A Figura~\ref{fig:projeto-time3} apresenta o time junto aos \textit{stakeholders} durante a primeira reunião do projeto.

O período de execução foi de 29 de julho de 2025 a 13 de novembro de 2025. Para fins de registro e identificação da equipe envolvida, esta seção é acompanhada por uma foto do time.


\begin{figure}[H]
  \centering
  \small
  \caption{Time do projeto WeConecta}
  \includegraphics[width=1\linewidth]{conteudo//4 - ages III//conteudo//figures/foto-time.png}
  Fonte: Autor
  \label{fig:projeto-time3}
\end{figure}
